\chapter{Introduction}
\label{ch:introduction}

Text generation with large language models has, in the space of a few years, moved from a research curiosity to a technology that underpins a wide range of applications, from writing assistants and conversational agents to summarization, translation, and code synthesis. The dominant paradigm behind these systems is \textbf{autoregressive generation}: a model factorizes the probability of a sequence into a product of next-token conditionals and produces text one token at a time, from left to right, sampling or maximizing each token given the tokens already committed \citep{radford2019language, brown2020language}. This factorization is what makes modern language models trainable at scale and fast at inference, and it is the reason they work as well as they do.

The same factorization, however, is the source of a persistent and practically important limitation. Because generation proceeds strictly from left to right and each token is fixed once emitted, an autoregressive model has no mechanism for revising an earlier decision in light of a later one. This is not a problem when the objective is simply to produce fluent continuations, but it becomes a problem the moment we want to \emph{control} what the model produces. Many of the properties we care about most, such as the overall sentiment of a passage, its topical focus, its adherence to a factual constraint, or the absence of toxic content, are \textbf{global properties} of a completed sequence rather than local properties of any single token. A greedy left-to-right decision about the next token cannot see these properties, because the object that would be evaluated against them, the finished text, does not yet exist. The problem of steering a language model toward such global objectives is the problem of \textbf{controllable text generation}, and it is the subject of this thesis.

There are, broadly, two established ways of imposing control, and each carries a cost. The first is to bake the desired behaviour into the parameters of the model through supervised fine-tuning or reinforcement learning from human feedback \citep{ouyang2022training}. This is effective, but it is \emph{static}: a model tuned to be helpful and non-toxic embodies exactly those constraints and no others, and a new constraint at inference time requires a new dataset and a new training run. The second is to filter or re-rank many independent generations, discarding those that violate the constraint, which is conceptually simple but computationally wasteful and, for constraints that are rarely satisfied by chance, effectively intractable. What both approaches lack is the ability to introduce a fresh constraint at the moment of generation, cheaply, without touching the underlying model.

\section{Energy-Based Control and Its Central Assumption}

An elegant alternative reframes decoding itself. Rather than generating a sequence token by token, one can define a probability distribution over the space of all sequences and \emph{sample} from it. Following the standard formulation of \textbf{energy-based models} \citep{lecun2006tutorial, hinton2002training}, the distribution is written in the Boltzmann form
\begin{equation}
    p(\bm{x}) = \frac{1}{Z}\exp\bigl(-\energy(\bm{x})\bigr),
    \label{eq:intro-ebm}
\end{equation}
where $\energy(\bm{x})$ is the \textbf{energy} of a sequence $\bm{x}$, low energy corresponds to a desirable sequence, and $Z$ is a normalizing constant that need never be computed if one only wishes to sample. The attraction of this view for controllable generation is that the energy can be assembled additively from independent terms:
\begin{equation}
    \energy(\bm{x}) = -\log p_{\text{LM}}(\bm{x}) + \lambda\, C(\bm{x}),
    \label{eq:intro-energy}
\end{equation}
where the first term, the negative log-likelihood of a frozen pretrained language model, supplies \emph{fluency}, and the second term $C(\bm{x})$ is an arbitrary differentiable \emph{constraint} weighted by $\lambda$. Because the constraint is simply added on, it can be specified at inference time and swapped freely: a practitioner who wants positive sentiment today and factual consistency tomorrow changes $C(\bm{x})$ and leaves the language model untouched. This is the promise of plug-and-play, inference-time controllable generation, and a substantial and growing body of recent work pursues it, notably through gradient-guided decoders such as PPLM \citep{dathathri2020plug}, FUDGE \citep{yang2021fudge}, GeDi \citep{krause2021gedi}, and, most directly related to this thesis, sampling-based methods such as COLD \citep{qin2022cold} and MuCoLa \citep{kumar2022gradient}, which draw samples from the energy in \eqref{eq:intro-energy} using continuous relaxations and Langevin-style dynamics.

Every method in this family rests on a single, rarely examined assumption. To sample from the energy in \eqref{eq:intro-energy} with a gradient-guided method, one must be able to compute a gradient of $-\log p_{\text{LM}}(\bm{x})$ that points, at least on average, toward better sequences, and one must be able to follow that gradient across the space of possible texts. In other words, the entire framework presupposes that the frozen likelihood of an autoregressive language model, viewed as a function over sequences, defines a \emph{landscape that a sampler can actually navigate}. This assumption is intuitive, and it is almost never tested directly. It is the central object of study in this thesis, and the central finding is that, for standard autoregressive language models on discrete text, the assumption does not hold.

\section{From Sampling to Amortization}

The natural response to a gradient that will not cooperate is to stop using the gradient. If a per-sequence energy can still be \emph{evaluated}, even when it cannot be usefully \emph{differentiated}, then one can treat the language model as a black box that scores finished sequences and learn a separate policy that generates in proportion to that score. This is the idea of \textbf{amortized inference}, and its most principled recent realization for sequence generation is the \textbf{Generative Flow Network}, or GFlowNet \citep{bengio2021flow, bengio2023gflownet}, which trains a generative policy so that the probability of producing a complete sequence is proportional to its reward. Applied to language models, GFlowNets have been proposed as a way to sample diverse, high-reward text without the mode-collapse of reward-maximizing reinforcement learning \citep{hu2024amortizing}. The second half of this thesis asks whether amortizing the search in this way escapes the difficulties that afflict the gradient-guided samplers, and finds that it does not, because the difficulty lies not in the search procedure but in the energy itself.

\section{Research Questions and Contributions}

The work in this thesis is organized around four research questions.

\begin{description}
    \item[RQ1.] Can the frozen likelihood of an autoregressive language model be sampled effectively with theoretically faithful Langevin dynamics on the discrete text manifold, and if not, why not?
    \item[RQ2.] What is the role of the Metropolis--Hastings correction in each of the discrete and continuous samplers, and does making the samplers theoretically correct also make them work?
    \item[RQ3.] Does amortizing the energy with a GFlowNet, which never differentiates the reward, escape the failures observed for the gradient-guided samplers?
    \item[RQ4.] Can an additive constraint term steer generation toward a target attribute on this energy landscape, as the plug-and-play framework requires?
\end{description}

In answering these questions the thesis makes the following contributions. It establishes, through a controlled ablation that separates the direction of the gradient from its magnitude, that the frozen autoregressive gradient carries no usable information for discrete sampling, and it replicates this finding across five energy functions. It identifies and measures the \emph{linearization failure} that is the geometric cause of this result, showing that the region over which the sampler's Taylor surrogate is valid is smaller than the smallest possible discrete step. It measures the \emph{likelihood trap} on the study's own models, demonstrating that the lowest-energy text is the most degenerate. It provides a precise, measured account of the \emph{Metropolis--Hastings breakdown} in continuous space, attributing the collapse of the acceptance ratio to the proposal term and connecting it to the non-Lipschitz drift induced by nearest-neighbour projection. It documents three mechanistically distinct failure modes of GFlowNet fine-tuning on a frozen-likelihood reward, and it shows in a unifying experiment that Langevin sampling on the amortized energy is indistinguishable from sampling on the untuned base. Finally, it reports a constrained-generation extension that confirms the constraint gradient carries no reliable steering signal, and it proposes a positive-control experiment on a diffusion language model that would falsify the thesis's central explanation.

\section{Structure of the Thesis}

Chapter~\ref{ch:background} introduces the background needed to follow the technical development: energy-based models, Langevin dynamics and the Metropolis--Hastings correction, the discrete and continuous samplers studied here, and GFlowNets. Chapter~\ref{ch:related} situates the work within the literature on controllable and energy-based text generation and reviews the specific difficulties that discreteness and the geometry of language-model embeddings pose for gradient-guided methods. Chapter~\ref{ch:methodology} describes the experimental platform, the masked-token recovery task, the five energy functions, the sampler configurations, the evaluation metrics, and the four diagnostic experiments designed to isolate the mechanisms. Chapter~\ref{ch:results} presents the results, beginning with the calibration of the samplers, proceeding through the annealing dynamics, the Metropolis--Hastings breakdown, and the central gradient-fallacy result, and then turning to the GFlowNet experiments and the constrained-generation extension. Chapter~\ref{ch:discussion} draws the mechanisms together into a single account and considers the implications for future work, including the positive-control experiment. Chapter~\ref{ch:conclusion} concludes.
